\textbf{C040 双曲线焦准比定义。} $a,b>0,c^2=a^2+b^2,e=c/a>1$。
横向双曲线 $x^2/a^2-y^2/b^2=1$，右焦点 $(c,0)$ 配右准线 $x=a^2/c$，左侧对应负号；同一焦准对适用于两支。
$PF_R=|ex-a|=e|x-a^2/c|$。由第一定义作和差可快速推出；反向用
$PF_R^2-e^2d(P,\ell_R)^2=-b^2(x^2/a^2-y^2/b^2-1)$ 验证。
一般判定须 $F\notin\ell$；若焦点在准线上且 $e>1$，距离等式退化为两条直线。
例：$x^2/4-y^2/5=1$ 上 $PF_R=6$，$x=16/3$ 或 $-8/3$，两者均有实际点实现。
